\section{SOLL-Konzept}

\subsection{Zielbild des Avisierungsprozesses}

Das Zielbild umfasst einen durchgängigen digitalen
Avisierungsprozess für bereits im DISPO-System geplante
Heizöllieferungen. DISPO bleibt dabei das führende System für
Auftrags-, Tour- und Lieferplanungsdaten. Die digitale
Lieferavisierung übernimmt die für die Kundenabstimmung
erforderlichen Informationen, steuert den
Rückbestätigungsprozess und übermittelt den daraus
resultierenden Status über die vorgesehene
DISPO-Rückmeldeschnittstelle.

Die Kommunikation mit dem Kunden soll standardisiert über
E-Mail, SMS oder WhatsApp erfolgen. Jede Bestätigungsanfrage
enthält einen individuellen Link, über den der Kunde die
relevanten Auftrags- und Lieferinformationen einsehen und den
vorgeschlagenen Liefertermin bestätigen oder ablehnen kann.
Der gewählte Kommunikationskanal beeinflusst dabei nicht den
fachlichen Ablauf der Kundenrückbestätigung.

Kundenreaktionen und das Ausbleiben einer Antwort innerhalb der
festgelegten Frist sollen in einem einheitlichen Statusmodell
verarbeitet werden. Über ein geschütztes Avisierungsdashboard
erhält die Disposition eine zentrale Übersicht über Touren,
Aufträge, aktuelle Avisierungsstatus und frühere
Bestätigungsanfragen. Für Vorgänge mit dem Status
\texttt{NO\_RESPONSE} oder einem technisch fehlgeschlagenen
Versand kann eine erneute Bestätigungsanfrage angestoßen werden.
Dabei kann der bisherige Kommunikationskanal beibehalten oder
ein anderer Kanal ausgewählt werden.

Für bestätigte Lieferungen sollen am Liefertag zusätzlich
Tracking-Informationen bereitgestellt werden. Dazu wird die über
die vorgesehene DISPO-Tracking-Schnittstelle abgerufene
Fahrerposition mit den aus der Lieferadresse ermittelten
Zielkoordinaten zusammengeführt. Auf diese Weise ergänzt die
Lieferverfolgung den Rückbestätigungsprozess um aktuelle, für
den Kunden zugängliche Informationen zum Lieferkontext.

Das Zielbild verbindet damit die standardisierte
Kundenkommunikation, die workflowgestützte Fristverarbeitung,
die Statusrückmeldung und die zentrale Übersicht für die
Disposition in einem gemeinsamen Avisierungsprozess. Die
Liefer- und Tourenplanung verbleibt weiterhin im DISPO-System
und bleibt fachlich von der avisierungsspezifischen Verarbeitung
abgegrenzt.

\subsection{SOLL-Prozess}

\subsubsection{Prozessablauf}

Ein bereits im DISPO-System geplanter Auftrag wird mit den für
die Avisierung erforderlichen Kunden-, Kontakt-, Tour- und
Lieferdaten an den Avisierungsservice übermittelt. Zusätzlich
werden der vorgesehene Kommunikationskanal und die gewünschte
Antwortdauer festgelegt. Auf dieser Grundlage wird eine
individuelle Bestätigungsanfrage erstellt.

Die Kundenbenachrichtigung wird anschließend über den für den
Auftrag festgelegten Kommunikationskanal versendet. Nach dem
erfolgreichen Versand beginnt die Antwortfrist und die
Bestätigungsanfrage steht für eine Kundenreaktion zur Verfügung.
Die Frist ergibt sich aus der festgelegten Antwortdauer und endet
spätestens mit dem Beginn des Lieferzeitfensters.

Über den personalisierten Link gelangt der Kunde zur
Bestätigungsseite. Dort werden die für seine Entscheidung
relevanten Auftrags- und Lieferinformationen angezeigt. Der
Kunde kann den vorgeschlagenen Liefertermin bestätigen oder
ablehnen und optional einen Kommentar übermitteln.

Eine Bestätigung führt zum Status \texttt{CONFIRMED}, eine
Ablehnung zum Status \texttt{REJECTED}. Geht innerhalb der
festgelegten Frist keine Antwort ein, wird der Vorgang als
\texttt{NO\_RESPONSE} abgeschlossen. Die Kundenreaktion und ein
gegebenenfalls übermittelter Kommentar werden gemeinsam mit dem
erreichten Status gespeichert.

Der Abschlussstatus wird anschließend zusammen mit der externen
Auftragsnummer über die vorgesehene Rückmeldeschnittstelle an
das DISPO-System übermittelt. Der aktuelle Bearbeitungsstand
sowie die bisherigen Bestätigungsanfragen und
Kundenreaktionen bleiben im Avisierungsdashboard für die
Disposition nachvollziehbar.

Nach einem Fristablauf ohne Kundenreaktion oder einem endgültig
fehlgeschlagenen Nachrichtenversand kann die Disposition über
das Avisierungsdashboard eine erneute Bestätigungsanfrage
auslösen. Dabei können der bisherige Kommunikationskanal und die
Antwortdauer beibehalten oder angepasst werden. Frühere Anfragen
bleiben in der Historie erhalten, während nur der Link der
aktuellen Bestätigungsanfrage für eine weitere Kundenreaktion
verwendet werden kann. Eine erneute Avisierung ist nur möglich,
solange das zugehörige Lieferzeitfenster noch nicht begonnen hat.

Nach der Bestätigung eines Liefertermins können am Liefertag
zusätzlich Tracking-Informationen bereitgestellt werden. Dazu
werden die aus der Lieferadresse ermittelte Zielposition und die
über die vorgesehene DISPO-Tracking-Schnittstelle bereitgestellte
Fahrerposition zusammengeführt und dem Kunden für die
Lieferverfolgung zugänglich gemacht.

\subsubsection{BPMN-Modell}

Der zuvor beschriebene SOLL-Prozess bildet die fachliche
Grundlage für den workflowgestützten Ablauf einer
Bestätigungsanfrage. Abbildung~\ref{fig:confirmation-request-process}
zeigt dessen technische Konkretisierung als ausführbares
BPMN-Modell.

Nach der Erstellung einer Bestätigungsanfrage steuert der
Workflow zunächst den Versand der Kundenbenachrichtigung.
Vorübergehend fehlgeschlagene Versandversuche können über einen
dafür vorgesehenen Wiederholungspfad erneut ausgeführt werden.
Kann die Nachricht endgültig nicht versendet werden, wird der
Vorgang über einen eigenen Fehlerpfad abgeschlossen.

Nach einem erfolgreichen Versand wartet ein ereignisbasiertes
Gateway auf den Eingang einer Kundenreaktion, den Ablauf der
Antwortfrist oder die Ablösung der Anfrage durch eine neuere
Bestätigungsanfrage. Eine Kundenreaktion oder ein Fristablauf
führt zur Ermittlung des fachlichen Abschlussstatus und
anschließend zur Rückmeldung an das DISPO-System. Eine abgelöste
Anfrage wird dagegen beendet, ohne einen veralteten Status an
DISPO zu übermitteln.

\begin{figure}[H]
    \centering
    \includegraphics[
        width=\textwidth,
        height=0.78\textheight,
        keepaspectratio
    ]{figures/confirmation-request-process.png}
    \caption{Ausführbarer BPMN-Prozess einer Bestätigungsanfrage}
    \label{fig:confirmation-request-process}
\end{figure}